\appendix

\section{Ethics Statement}

This research on memory-augmented learning for large language model agents raises several important ethical considerations that we wish to acknowledge.

Though our suggestibility work was focused on how the model's instruction-following ability varied with dataset, this kind of approach could also be used to more efficiently jailbreak models to spread misinformation. Future work should be careful to avoid developing tools to improve the suggestibility of models to the point that they spread harmful misinformation.

We also recognize that improved adaptation capabilities may exacerbate existing biases in these agents. Because the insights are generated by the agent itself, even with feedback from the labeled data, it could cause the agent to reinforce its preconceptions about the world, which may perpetuate harmful stereotypes. Future work should explore safeguards to identify and mitigate such bias amplification.

\section{Implementation Details}
\subsection{RAG Implementation}
\label{app:rag_details}
We used blevlabs/stella\_en\_v5 as our encoder model and FAISS as our vector database. Similarity was based purely on the encodings of the questions in each dataset.

Though we did experiment with fine-tuning an encoder model to increase the separation of the classes in each dataset (for example, preferred vs not-preferred games for each Steam user) in embedding space, we did not see significant improvements in performance.

\subsection{Dataset: Additional Details}
\label{app:data}

\paragraph{Multi-Condition Ranking:} The original data source was presented in a free-form response format. However, for this task, we transformed it into a multiple-choice format for easier evaluation. The model now has to choose the correct answer from four unique orderings of the items.

\paragraph{NFCorpus:} The original NFCorpus data source associated each article with many papers with varying degrees of separation, which we transformed into this pairwise setup by choosing one paper at the closest and furthest level of separation possible for each paper. Sampled 500 shortest combinations of articles and papers to avoid context-length issues.

\section{Prompts}


\begin{promptbox}{Critique Generation} 
\small

\textbf{Question:} \{Question\}

\textbf{Initial Agent Response:} \{PA Initial Prediction\}

The correct answer is \{Ground Truth Answer\}. Explain why this is the correct answer, following the following JSON format\: 

\{

\quad correct\_answer: correct\_answer, 

\quad local\_reason: Specific reasons why this answer is correct in this particular case., 

\quad global\_reason: General reasons why this answer is correct that can be applied to other questions.

\}. 

Respond only with JSON.
\end{promptbox}

\begin{promptbox}{Semantic Memory Generation} 
\small
Your job is to summarize a set of self-critiques made by some agent as they perform different instances of their task.  For each instance you will be shown the output of the agent, followed by the critiques made by the agent after they were told the correct answer.  Distill those critiques into a helpful summary of advice to the agent, paying particular attention to instances where the agent outputs an incorrect answer. Produce your output in a form that can be used directly as instructions to the agent. You should summarize the key points in these critiques.  Be precise and concise.  Do not repeat yourself.

For example in train\_set:

\quad \{Question\} \{Answer\} \{Critique\} 

\end{promptbox}


\begin{promptbox}{Performance Agent with Semantic Memory} 
\small

\{Question\}

Here is some helpful advice that will help you make your decision: \{Summary\}
\end{promptbox}

\begin{promptbox}{Performance Agent with Episodic Memory} 
\small
For example in examples:

    \quad \textbf{Question:} \{Example Question\}
    
    \quad \textbf{Initial Agent Response:} \{PA Initial Prediction\}

    \quad (If \EPLABEL{}:) \textbf{Correct Answer:} \{correct\_answer\}

    \quad (If \EPCRIT{}:) \textbf{Reflection:} \{reflection\}
    
    
Here is your final question, make sure to learn from your past mistakes! \{Question\}
\end{promptbox}

\begin{promptbox}{Performance Agent with Episodic and Semantic Memory} 
\small
For example in examples:

    \quad \textbf{Question:} \{Example Question\}
    
    \quad \textbf{Initial Agent Response:} \{PA Initial Prediction\}

    \quad \textbf{Reflection:} \{reflection\}
    
Here is some additional advice to guide your response: \{Summary\} 

Here is your final question, make sure to learn from your past mistakes! \{Question\}

\end{promptbox}

\section{Examples}
\label{appendix_examples}


\begin{promptbox}{Example Semantic Memory

qwen3-vl-235b-a22b-thinking

Multi-Condition Ranking}
\begin{enumerate}[label=\textbf{\arabic*.}]

  \item \textbf{Interpret positional constraints correctly}
    \begin{itemize}
      \item ``Last from right'' means leftmost position in left-to-right list
      \item ``Last from left'' means rightmost position in left-to-right list
      \item These directional constraints override general priority ordering
    \end{itemize}

  \item \textbf{Categorize items by priority first}
    \begin{itemize}
      \item Identify high priority items (often with fixed position requirements)
      \item Identify medium priority items (typically based on specific attributes like ``related to X'' or ``born after Y'')
      \item Low priority items are all remaining items not classified as high/medium
    \end{itemize}

  \item \textbf{Apply sorting rules within priority groups}
    \begin{itemize}
      \item Low priority: sort by character count ascending (smallest to largest)
      \item Medium priority: sort by physical size (smallest to largest) or birthday (oldest to newest) as specified
      \item High priority: typically have fixed positions rather than being sorted among themselves
    \end{itemize}

  \item \textbf{Resolve ties systematically}
    \begin{itemize}
      \item When items have equal sorting values (e.g., same character count), use alphabetical order as tiebreaker
      \item Never assume equal values should maintain original order without explicit instruction
    \end{itemize}

  \item \textbf{Sequence priority groups correctly}
    \begin{itemize}
      \item The order of priority groups (high/medium/low) depends on specific problem constraints
      \item Verify which priority group should come first/last based on the conditions
      \item Fixed position constraints take precedence over general priority ordering
    \end{itemize}

  \item \textbf{Handle vacuous conditions properly}
    \begin{itemize}
      \item If no items meet a condition (e.g., no medium priority items), that condition doesn't affect the order
      \item Don't apply sorting rules to empty priority groups
    \end{itemize}

  \item \textbf{Verify all constraints sequentially}
    \begin{itemize}
      \item First satisfy fixed position requirements
      \item Then apply sorting rules to remaining items
      \item Always check that the final order satisfies all conditions simultaneously
    \end{itemize}

  \item \textbf{Prioritize explicit constraints over assumptions}
    \begin{itemize}
      \item When conditions conflict with your interpretation, re-examine the problem statement
      \item The correct answer must satisfy all explicit constraints, even if it contradicts your initial understanding of priority hierarchy
    \end{itemize}

\end{enumerate}
\end{promptbox}

\begin{promptbox}{Example Episodic Memory

qwen3-vl-235b-a22b-thinking

Multi-Condition Ranking}
\subsection*{Example 1:}
\textbf{Question:} Given the following conditions: ['Sort the items with a medium priority based on their height from the shortest to the tallest', 'Item ``bookshelf'' with a high priority should be the last from left', 'Sort the items with a low priority based on their character count from the smallest to largest'], sort the list of items ['bookshelf', 'water drop', 'chair', 'apartment', 'bottle'] from left to right. Choose the correct order from the provided options. Do not provide any explanation.\\
Option 1: ['bookshelf', 'apartment', 'water drop', 'bottle', 'chair']\\
Option 2: ['apartment', 'water drop', 'bottle', 'bookshelf', 'chair']\\
Option 3: ['water drop', 'bottle', 'chair', 'apartment', 'bookshelf']\\
Option 4: ['apartment', 'water drop', 'bottle', 'chair', 'bookshelf']\\
Respond only in JSON in the format \{``choice'': N\}, where N is 1 for Option 1, 2 for Option 2, and so on.

\textbf{Agent Response:} \{``choice'': 4\}

\textbf{Reflection:}
\begin{itemize}
  \item \textbf{Correct Answer:} 3
  \item \textbf{Local Reason:} The non-bookshelf items are sorted by height from shortest to tallest (water drop $<$ bottle $<$ chair $<$ apartment), satisfying the medium priority condition. `bookshelf' is correctly placed last as high priority. The low priority condition does not apply here as all non-high priority items are medium priority, making option 3 the only valid arrangement.
  \item \textbf{Global Reason:} When sorting with multiple priority levels, high priority items are fixed in their required position first. Remaining items are grouped by their priority level and sorted according to their respective criteria. The correct order combines these sorted groups while maintaining priority hierarchy, ensuring all specified conditions are met without contradiction.
\end{itemize}

\subsection*{Example 2:}
\textbf{Question:} Given the following conditions: ['Sort the items with a low priority based on their character count from the smallest to largest', 'Item ``bookshelf'' with a high priority should be the last from left', 'Sort the items with a medium priority based on their height from the shortest to the tallest'], sort the list of items ['ant', 'bookshelf', 'shoe', 'apartment', 'sofa'] from left to right. Choose the correct order from the provided options. Do not provide any explanation.\\
Option 1: ['bookshelf', 'shoe', 'sofa', 'apartment', 'ant']\\
Option 2: ['ant', 'shoe', 'sofa', 'apartment', 'bookshelf']\\
Option 3: ['sofa', 'ant', 'apartment', 'bookshelf', 'shoe']\\
Option 4: ['apartment', 'bookshelf', 'ant', 'sofa', 'shoe']\\
Respond only in JSON in the format \{``choice'': N\}, where N is 1 for Option 1, 2 for Option 2, and so on.

\textbf{Agent Response:} \{``choice'': 2\}

\textbf{Reflection:}
\begin{itemize}
  \item \textbf{Correct Answer:} 2
  \item \textbf{Local Reason:} The list must end with `bookshelf' (high priority). The remaining items are sorted such that low priority items (`ant', `shoe', `sofa') are ordered by character count (3, 4, 4), and medium priority (`apartment') follows. Option 2 correctly places `bookshelf' last and orders the others by character count for low priority items.
  \item \textbf{Global Reason:} When sorting with multiple priority levels, high priority items are fixed first (e.g., last position here), then lower priority groups are sorted by their specified criteria (e.g., character count or height). Correct answers must satisfy all constraints simultaneously, with higher priority rules taking precedence.
\end{itemize}

\subsection*{Example 3:}
...

\end{promptbox}



\section{Additional Results}

\subsection{Detailed Cost Analysis}

\begin{table}[t]
\centering
\scriptsize
\setlength{\tabcolsep}{3pt}
\begin{tabular}{lccccccc}
\toprule
\makecell{Model and \\ Metric} & \makecell{multi\\cond\\ranking} & \makecell{NFCorpus} & \makecell{PubMed} & \makecell{Steam\\Pref} & \makecell{Book\\Pref} & \makecell{Anime\\Pref} & \makecell{Movie\\Pref} \\
\midrule
\multicolumn{8}{l}{\textbf{gpt-5-2025-08-07}} \\
\quad Init Preds & 609.6k & 105.2k & 117.0k & 202.9k & 160.0k & 135.2k & 135.6k \\
\quad Critiques & 500.4k & 187.8k & 217.1k & 330.5k & 296.1k & 305.1k & 322.7k \\
\quad Summary & 3.3k & 3.7k & 5.4k & 29.0k & 27.3k & 31.7k & 33.6k \\
\quad BE: EP\_CRIT & 759 & 2.9k & never & 2.8k & 977 & 1.7k & 2.4k \\
\quad BE: SEM\_CRIT & 800 & never & never & 1.4k & 884 & 1.1k & 1.3k \\
\quad BE: EP+SEM\_CRIT & 620 & 2.8k & never & 1.9k & 866 & 1.4k & 2.2k \\
\hline
\multicolumn{8}{l}{\textbf{gpt-oss-20b}} \\
\quad Init Preds & 818.7k & 112.8k & 56.7k & 53.3k & 83.2k & 55.5k & 45.0k \\
\quad Critiques & 737.4k & 130.2k & 112.8k & 169.5k & 151.4k & 147.1k & 148.2k \\
\quad Summary & 1.5k & 893 & 2.1k & 14.3k & 15.5k & 14.3k & 13.7k \\
\quad BE: EP\_CRIT & 995 & 1.6k & 4.5k & 219 & 201 & 117 & 107 \\
\quad BE: SEM\_CRIT & 493 & never & never & 158 & 139 & 92 & 90 \\
\quad BE: EP+SEM\_CRIT & 845 & 956 & 6.2k & 225 & 256 & 135 & 112 \\
\hline
\multicolumn{8}{l}{\textbf{qwen3-vl-235b-a22b-thinking}} \\
\quad Init Preds & 1.30M & 195.8k & 121.8k & 211.3k & 278.2k & 230.1k & 238.9k \\
\quad Critiques & 1.08M & 219.1k & 147.1k & 179.5k & 187.0k & 199.1k & 141.5k \\
\quad Summary & 1.2k & 858 & 1.4k & 57.7k & 42.4k & 46.6k & 39.4k \\
\quad BE: EP\_CRIT & 2.0k & 842 & 5.0k & 1.6k & 967 & 958 & 781 \\
\quad BE: SEM\_CRIT & 2.4k & 4.2k & 1.1k & 564 & 719 & 436 & 370 \\
\quad BE: EP+SEM\_CRIT & 1.7k & 880 & 84.6k & 3.0k & 1.3k & 935 & 871 \\
\bottomrule
\end{tabular}
\caption{Token usage (completion tokens) per training stage and break-even test-question counts versus EP\_LABEL. Init Preds / Critiques / Summary are total train-set costs. All preference data numbers are summed across users, so all columns represent the cost across a training size of 250 questions. Semantic memory costs will be proportionally higher for preference data because it has to be generated separately for all ten users.}
\label{tab:usage}
\end{table}

\Cref{tab:usage} reports the total completion tokens used during offline memory generation (initial predictions, critiques, and semantic summary), along with the number of test queries at which each critique strategy amortizes its offline cost relative to \EPLABEL{}. Semantic summary generation is extremely cheap in all cases because it requires only a single LLM call (per-user) over the already-generated critiques; the dominant upfront costs are initial predictions and critique generation. This means that the accuracy gains we demonstrate for \EPSEMCRIT{} over \EPCRIT{} in the main paper come at essentially no additional memory-generation cost.

Break-even speed varies substantially by dataset type. Preference datasets break even quickly, sometimes within double-digit numbers of queries, while knowledge-intensive datasets require more queries to amortize, with PubMed being the highest by a significant margin. Notably, despite Multi-Condition Ranking being by far the most expensive to generate memories for, it demonstrates such significant efficiency gains at inference time that it breaks even within a few hundred queries. The ``never'' entries --- where a strategy does not break even --- occur in a small number of model-dataset-strategy combinations, primarily \SEMCRIT{} on NFCorpus and PubMed. In these cases the critique method uses more inference tokens than \EPLABEL{}, so the upfront generation cost is never amortized through token savings. However, we still observe significant accuracy gains over \EPLABEL{} on these datasets (for example, on NFCorpus \EPSEMCRIT{} improves over \EPLABEL{} for all six models), indicating that the accuracy and efficiency benefits of critique memory are complementary but do not always co-occur. Overall, for the majority of model-dataset combinations, the offline generation cost is amortized within a modest number of queries, making the approach practical for most production deployments.

\subsection{Training Data Size Ablation}

\begin{table}[t]
\centering
\scriptsize
\setlength{\tabcolsep}{3pt}
\begin{tabular}{lcccc | cccc}
\toprule
\makecell{Model and \\ Experiment} & \multicolumn{4}{c|}{\makecell{Multi-\\Cond.\\Ranking}} & \multicolumn{4}{c}{\makecell{Steam\\Pref}} \\
 & 25\% & 50\% & 75\% & 100\% & 25\% & 50\% & 75\% & 100\% \\
\midrule
\multicolumn{9}{l}{\textbf{gpt-4o-mini-2024-07-18}} \\
\quad \EPCRIT{} & 34.8 & \textbf{35.6} & 35.6 & 34.8 & 67.2 & 65.6 & 69.2 & \textbf{71.2} \\
\quad \EPSEMCRIT{} & 35.2 & \textbf{36.8} & 34.4 & 36.8 & 69.2 & 68.4 & 71.2 & \textbf{73.6} \\
\hline
\multicolumn{9}{l}{\textbf{gpt-5-2025-08-07}} \\
\quad \EPCRIT{} & \textbf{99.2} & 97.2 & 98.0 & 96.8 & 66.0 & 62.8 & 66.4 & \textbf{69.2} \\
\quad \EPSEMCRIT{} & 98.8 & 98.4 & 98.8 & \textbf{99.2} & 66.4 & 64.4 & 66.4 & \textbf{68.8} \\
\hline
\multicolumn{9}{l}{\textbf{gpt-oss-20b}} \\
\quad \EPCRIT{} & 52.8 & 52.4 & 54.0 & \textbf{57.6} & 59.2 & 57.6 & 60.4 & \textbf{63.2} \\
\quad \EPSEMCRIT{} & 57.6 & \textbf{58.8} & 55.2 & 58.8 & 59.6 & 58.0 & 60.8 & \textbf{62.8} \\
\bottomrule
\end{tabular}
\caption{Train-size ablation showing accuracy  of \EPCRIT{} and \EPSEMCRIT{} at different training-set fractions. \textbf{Bold} marks the best fraction per (model, dataset, experiment).}
\label{tab:train_size_ablation}
\end{table}

Our approach is naturally suited to online deployment. Critique generation for each training example is fully independent, so new episodic memories can be generated and added to the vector database in parallel as new labeled data arrives without regenerating existing memories, and semantic memory needs only periodic regeneration as the critique pool grows. To characterize how performance scales as labeled data accumulates, we evaluated \EPCRIT{} and \EPSEMCRIT{} at 25\%, 50\%, 75\%, and 100\% of the training data on a subset of models and datasets, where each larger split is a strict superset of the smaller ones. Because episodic memory is stateless and semantic memory is regenerated from the full critique pool at each stage, this is equivalent to an online agent that accumulates labeled examples over time and periodically refreshes its memory in batches.


As shown in \Cref{tab:train_size_ablation}, both \EPCRIT{} and \EPSEMCRIT{} show clear, largely monotonic improvement on Steam Preference as data accumulates across all three models, consistent with what we would expect from an online system that improves with experience. Multi-Condition Ranking shows a general upward trend but with more variance across splits. We hypothesize that embedding-based retrieval is less reliable on this dataset: Multi-Condition Ranking requires reasoning about relational properties between items rather than intrinsic item attributes, which may mean that the most useful critiques are not the most embedding-similar. As the pool grows, more superficially similar but less informative examples could crowd out the most useful ones. While the embedding-retrieved episodic memories remain clearly helpful over zero-shot, this suggests that alternative retrieval strategies such as diversity-aware or task-adapted retrieval are a promising direction for future work on tasks with this property. A second notable pattern is that on Multi-Condition Ranking, \EPSEMCRIT{} at 25\% of the training data always matches or exceeds full-data \EPCRIT{}, suggesting that semantic memory provides useful task-level signals even from small critique pools --- a property that would be especially valuable in early-stage online deployment when few labeled examples are available.

\subsection{Critique Component Ablation}

\begin{table}[t]
\centering
\scriptsize
\setlength{\tabcolsep}{3pt}
\begin{tabular}{lccccccc}
\toprule
\makecell{Model and \\ Variant} & \makecell{Multi-\\Cond.\\Ranking} & \makecell{NFCorpus} & \makecell{PubMed} & \makecell{Steam\\Pref} & \makecell{Book\\Pref} & \makecell{Anime\\Pref} & \makecell{Movie\\Pref} \\
\midrule
\multicolumn{8}{l}{\textbf{gpt-4o-mini-2024-07-18}} \\
\quad Local & \textbf{35.6} & \textbf{89.6} & 64.8 & \textbf{71.6} & 64.8 & 60.4 & \textbf{59.6} \\
\quad Global & 34.4 & 86.0 & \textbf{65.2} & 64.4 & 60.4 & 58.4 & 50.4 \\
\quad Full & 33.6 & 88.8 & 63.2 & \textbf{71.6} & \textbf{66.4} & \textbf{62.8} & 58.8 \\
\hline
\multicolumn{8}{l}{\textbf{gpt-5-2025-08-07}} \\
\quad Local & 97.2 & 97.6 & \textbf{72.8} & 68.8 & \textbf{70.0} & \textbf{75.6} & 64.4 \\
\quad Global & 97.6 & 97.6 & 68.8 & 66.0 & 67.6 & 67.6 & 63.2 \\
\quad Full & \textbf{98.0} & \textbf{98.4} & 70.8 & \textbf{69.2} & 69.2 & 71.6 & \textbf{67.2} \\
\hline
\multicolumn{8}{l}{\textbf{gpt-oss-20b}} \\
\quad Local & \textbf{59.2} & 94.8 & 66.4 & 61.6 & \textbf{65.2} & \textbf{64.0} & \textbf{66.4} \\
\quad Global & 44.4 & 94.8 & 66.0 & 63.6 & 60.0 & 60.0 & 58.4 \\
\quad Full & 56.4 & \textbf{95.2} & \textbf{67.2} & \textbf{64.8} & 60.4 & 60.8 & 63.6 \\
\bottomrule
\end{tabular}
\caption{Critique-parts ablation showing \EPCRIT{} accuracy when critiques contain only the local reason, only the global reason, or both (full). \textbf{Bold} marks the best variant per (model, dataset).}
\label{tab:critique_parts_ablation}
\end{table}

\Cref{tab:critique_parts_ablation} reports the accuracy of \EPCRIT{} when the critique attached to each retrieved example contains only the local reason, only the global reason, or the full critique (note, full is an independent trial from the accuracies reported in \Cref{tab:model_performance} and \Cref{tab:model_performance_reasoning}). A clear, model-dependent pattern emerges in how the two components are utilized. GPT-4o-mini and GPT-oss tend to benefit most from local rationales: the local-only variant is the best or tied-for-best on four of seven datasets for each, and the gap to global-only can be large --- nearly 15 points on Multi-Condition Ranking for GPT-oss (59.2 vs.\ 44.4) and roughly 9 points on Movie Preference for GPT-4o-mini (59.6 vs.\ 50.4). The stronger GPT-5 shows the opposite tilt: the full critique is best on four of seven datasets, and the gap between local-only and global-only is much smaller. This is consistent with a more capable model being better positioned to extract value from task-level patterns and to integrate them with per-instance rationales rather than relying on either signal alone.

Two broader observations follow. First, the global-only variant wins outright in only one of the 21 (model, dataset) cells (GPT-4o-mini on PubMed), indicating that instance-specific reasoning is the most consistently load-bearing component across the models and tasks we evaluate. Second, despite this, global reflections contribute meaningful complementary information once paired with local context --- most strikingly for GPT-5, but also for GPT-4o-mini on Book and Anime Preference and for GPT-oss on NFCorpus, PubMed, and Steam Preference, where the full variant matches or surpasses the local-only variant. Combining both components therefore provides broader coverage across the settings we evaluate than committing to a single signal type.


\subsection{Varying K}
\label{varying_k}


\begin{table}[htbp]
\scriptsize
\setlength{\tabcolsep}{3pt}
\centering
\begin{tabular}{lccccccc}
\toprule
\makecell{Model and \\ Experiment} & \makecell{Multi-\\Cond.\\Ranking} & \makecell{NFCorpus} & \makecell{PubMed} & \makecell{Steam\\Pref} & \makecell{Book\\Pref} & \makecell{Anime\\Pref} & \makecell{Movie\\Pref} \\
\midrule
\multicolumn{8}{l}{\EPCRIT{}} \\
\quad $K=1$ & 35.2 & 82.4 & 63.2 & 65.2 & 62.0 & 59.2 & 57.2 \\
\quad $K=3$ & 35.2 & 86.0 & 62.8 & 69.6 & 63.2 & 61.2 & 56.8 \\
\quad $K=5$ & 34.8 & 89.6 & 63.2 & {71.2} & 64.8 & 62.0 & {58.4} \\
\quad $K=10$ & {37.6} & {90.4} & {66.4} & 70.4 & {66.8} & {64.0} & 57.6 \\
\bottomrule
\end{tabular}
\caption{Effect on accuracy of varying $K$, the number of examples used in episodic memory, with gpt-4o-mini as the base model.}
\label{tab:varying_k}
\end{table}


In \Cref{tab:varying_k}, we observe that varying $K$, the number of examples included from the episodic memory module, can have significant effects on the accuracy of the overall performance agent, with an up to 8\% absolute difference in accuracy between different K values for the same setup. $K=5$ and $K=10$ both have similar results, and are on average better than $K=1$ or $K=3$. Because the results for $K=10$ were not significantly higher than $K=5$, we elected to use the simpler method and use $K=5$ for all episodic and episodic+semantic memory experiments in this paper.



\subsection{Model Token Timeout Rates}

\begin{table}[t]
\centering
\scriptsize
\setlength{\tabcolsep}{3pt}
\begin{tabular}{lccccccc}
\toprule
\makecell{Model and \\ Experiment} & \makecell{Multi-\\Cond.\\Ranking} & \makecell{NFCorpus} & \makecell{PubMed} & \makecell{Steam\\Pref} & \makecell{Book\\Pref} & \makecell{Anime\\Pref} & \makecell{Movie\\Pref} \\
\midrule
\multicolumn{8}{l}{\textbf{gpt-5-2025-08-07}} \\
\quad \textit{Train} & 1.2\ & 0.0\ & 0.0\ & 0.0\ & 0.0\ & 0.0\ & 0.0\ \\
\quad \ZEROSHOT{} & 2.0 & 0.0 & 0.0 & 0.0 & 0.0 & 0.0 & 0.0 \\
\quad \EPLABEL & 4.0 & 0.0 & 0.0 & 0.0 & 0.0 & 0.0 & 0.0 \\
\quad \EPCRIT{} & 1.2 & 0.0 & 0.0 & 0.0 & 0.0 & 0.0 & 0.0 \\
\quad \SEMCRIT{} & 0.4 & 0.0 & 0.0 & 0.0 & 0.0 & 0.0 & 0.0 \\
\quad \EPSEMCRIT{} & 0.4 & 0.0 & 0.0 & 0.0 & 0.0 & 0.0 & 0.0 \\
\hline
\multicolumn{8}{l}{\textbf{gpt-oss-20b}} \\
\quad \textit{Train} & 9.5\ & 0.5\ & 0.0\ & 0.1\ & 0.1\ & 0.0\ & 0.0\ \\
\quad \ZEROSHOT{} & 14.4 & 0.8 & 0.0 & 0.4 & 0.0 & 0.0 & 0.0 \\
\quad \EPLABEL & 60.0 & 1.2 & 0.0 & 6.8 & 8.4 & 12.0 & 13.6 \\
\quad \EPCRIT{} & 36.8 & 0.4 & 0.4 & 3.2 & 4.4 & 2.0 & 0.4 \\
\quad \SEMCRIT{} & 12.0 & 1.2 & 0.0 & 0.0 & 0.8 & 0.0 & 0.0 \\
\quad \EPSEMCRIT{} & 35.2 & 0.0 & 0.0 & 1.2 & 3.6 & 3.2 & 0.4 \\
\hline
\multicolumn{8}{l}{\textbf{qwen3-vl-235b-a22b-thinking}} \\
\quad \textit{Train} & 8.7\ & 0.1\ & 0.0\ & 0.0\ & 0.0\ & 0.0\ & 0.0\ \\
\quad \ZEROSHOT{} & 25.1 & 0.0 & 0.0 & 0.0 & 0.0 & 0.0 & 0.0 \\
\quad \EPLABEL & 40.5 & 1.2 & 0.0 & 0.0 & 0.0 & 0.4 & 0.0 \\
\quad \EPCRIT{} & 26.2 & 0.0 & 0.0 & 0.0 & 0.0 & 0.0 & 0.0 \\
\quad \SEMCRIT{} & 29.6 & 1.2 & 0.0 & 0.0 & 0.0 & 0.0 & 0.0 \\
\quad \EPSEMCRIT{} & 22.2 & 0.0 & 0.0 & 0.0 & 0.0 & 0.0 & 0.0 \\
\bottomrule
\end{tabular}
\caption{Timeout rate -  the percentage of responses where output tokens reached the token limit - for thinking models across datasets and experiments.}
\label{tab:timeout_rate_thinking}
\end{table}

In \Cref{tab:timeout_rate_thinking} we report token timeout rates by model and dataset. We include only reasoning models, as none of the non-reasoning models timed out on any question. The token limit was set to 8,192 tokens, counting both the internal reasoning trace and the final output.

Timeout rates vary substantially across both models and datasets. All three models timed out most frequently on Multi-Condition Ranking, by a wide margin. This is likely because these questions demand the greatest amount of explicit reasoning: rather than relying on external knowledge, they require the model to execute a sequence of logical steps to sort items according to multiple criteria. Given this structure, it is unsurprising that this dataset exhibits the highest timeout rates.

Even on the multi-condition ranking task, GPT-5 almost never timed out, suggesting substantially more efficient use of reasoning tokens. In contrast, GPT-oss experienced more timeouts than Qwen3-Thinking on the preference datasets, despite using far fewer tokens per question on average (see \Cref{tab:token_usage_thinking}). This apparent contradiction is explained by the much higher variance in GPT-oss’s response lengths. While GPT-oss often produces relatively concise outputs, it is also more prone to generating unusually long reasoning traces that exceed the token limit. On manual inspection, GPT-oss appeared to be particularly prone to getting stuck in loops during its reasoning trace.



\subsection{Per-User Preference Data Results}

When presenting the results for the preference datasets elsewhere in the paper, the results are averaged across 10 users per dataset. This is done because different users have very different preference patterns, and some users are more consistent than others. Sampling across many users gives us a better look at the models true performance at this task for the general population.
In \Cref{tab:per_user_steam}, \Cref{tab:per_user_book}, \Cref{tab:per_user_anime}, and \Cref{tab:per_user_movie} we present the results separated out by user, along with the average and standard deviation across all ten users.

\begin{table*}[ht]
\centering
\begin{subtable}[t]{0.49\textwidth}
\centering
\scriptsize
\setlength{\tabcolsep}{2.5pt}
\begin{tabular}{lcccccccccccc}
\toprule
\makecell{Model and \\ Experiment} & U.1 & U.2 & U.3 & U.4 & U.5 & U.6 & U.7 & U.8 & U.9 & U.10 & Avg & Std \\
\midrule
\multicolumn{13}{l}{\textbf{gpt-4o-mini-2024-07-18}} \\
\quad \ZEROSHOT{} & 52 & 68 & \textbf{80} & 40 & \underline{64} & 72 & 24 & 48 & 64 & 72 & 58 & 16 \\
\quad \EPLABEL & 56 & 68 & 60 & 48 & 52 & 64 & 36 & \textbf{80} & 68 & \underline{80} & 61 & 13 \\
\quad \EPCRIT{} & \underline{76} & \underline{80} & \underline{80} & \textbf{60} & \textbf{68} & 76 & \textbf{64} & 60 & \textbf{76} & 72 & 71 & 7 \\
\quad \SEMCRIT{} & 72 & 68 & 68 & 48 & 56 & \textbf{80} & 36 & 48 & 52 & 76 & 60 & 14 \\
\quad \EPSEMCRIT{} & \textbf{80} & \textbf{84} & 76 & \underline{60} & 64 & \underline{80} & \underline{60} & \underline{72} & \underline{72} & \textbf{88} & 74 & 9 \\
\hline
\multicolumn{13}{l}{\textbf{gpt-5-2025-08-07}} \\
\quad \ZEROSHOT{} & 56 & \underline{60} & 60 & 56 & 44 & 60 & 44 & 56 & 52 & 48 & 54 & 6 \\
\quad \EPLABEL & 60 & 44 & \underline{80} & 64 & 56 & \textbf{68} & \textbf{80} & 60 & \underline{64} & 68 & 64 & 10 \\
\quad \EPCRIT{} & \textbf{80} & \textbf{72} & 76 & \textbf{72} & 60 & \underline{64} & \underline{76} & 64 & 56 & \textbf{72} & 69 & 7 \\
\quad \SEMCRIT{} & 64 & 60 & 68 & 60 & \underline{64} & 64 & 48 & \textbf{76} & \textbf{72} & 68 & 64 & 7 \\
\quad \EPSEMCRIT{} & \underline{72} & 60 & \textbf{84} & \underline{68} & \textbf{72} & 56 & 76 & \underline{68} & 60 & \underline{72} & 69 & 8 \\
\hline
\multicolumn{13}{l}{\textbf{gpt-oss-20b}} \\
\quad \ZEROSHOT{} & 52 & 52 & 44 & \underline{64} & 40 & 64 & 36 & \underline{52} & \textbf{72} & 48 & 52 & 11 \\
\quad \EPLABEL & 56 & 52 & 60 & 60 & 44 & 56 & \textbf{64} & \textbf{60} & 44 & \textbf{76} & 57 & 9 \\
\quad \EPCRIT{} & \textbf{76} & \textbf{72} & \textbf{68} & 56 & \underline{60} & \underline{76} & 56 & 44 & 56 & \underline{68} & 63 & 10 \\
\quad \SEMCRIT{} & 56 & 56 & 48 & \textbf{76} & 36 & \textbf{80} & 52 & 44 & \underline{68} & 56 & 57 & 13 \\
\quad \EPSEMCRIT{} & \underline{68} & \underline{68} & \underline{64} & 56 & \textbf{64} & 76 & \underline{60} & 52 & 52 & 68 & 63 & 7 \\
\hline
\multicolumn{13}{l}{\textbf{llama-4-scout}} \\
\quad \ZEROSHOT{} & 60 & 72 & \textbf{72} & 48 & \underline{48} & 56 & 32 & 44 & 56 & \underline{76} & 56 & 13 \\
\quad \EPLABEL & \textbf{72} & \textbf{76} & 60 & \textbf{64} & \textbf{52} & \textbf{80} & \textbf{44} & \textbf{68} & \textbf{72} & \textbf{80} & 67 & 11 \\
\quad \EPCRIT{} & 68 & 72 & 60 & \underline{64} & 48 & 76 & 36 & 56 & 68 & 76 & 62 & 12 \\
\quad \SEMCRIT{} & \underline{72} & 72 & 60 & 52 & 48 & \underline{80} & \underline{44} & \underline{64} & 68 & 72 & 63 & 11 \\
\quad \EPSEMCRIT{} & 72 & \underline{76} & \underline{64} & 60 & 40 & 80 & 40 & 60 & \underline{72} & 72 & 64 & 13 \\
\hline
\multicolumn{13}{l}{\textbf{qwen3-235b-a22b-instruct-2507}} \\
\quad \ZEROSHOT{} & 52 & 48 & 68 & 44 & 56 & \underline{72} & 32 & \textbf{64} & \underline{60} & 52 & 55 & 11 \\
\quad \EPLABEL & 60 & 44 & 64 & 44 & 40 & 52 & 52 & \underline{64} & \textbf{72} & 64 & 56 & 10 \\
\quad \EPCRIT{} & \textbf{80} & 60 & \textbf{88} & \textbf{68} & \textbf{60} & 64 & \underline{72} & 60 & 60 & 72 & 68 & 9 \\
\quad \SEMCRIT{} & 52 & \textbf{72} & \underline{84} & 52 & \underline{60} & \textbf{80} & 56 & 52 & 56 & \textbf{76} & 64 & 12 \\
\quad \EPSEMCRIT{} & \underline{80} & \underline{64} & 76 & \underline{68} & 60 & 64 & \textbf{80} & 56 & 60 & \underline{76} & 68 & 8 \\
\hline
\multicolumn{13}{l}{\textbf{qwen3-vl-235b-a22b-thinking}} \\
\quad \ZEROSHOT{} & \textbf{68} & \textbf{80} & 60 & \underline{56} & 56 & \underline{60} & 32 & 48 & \textbf{68} & 60 & 59 & 12 \\
\quad \EPLABEL & 48 & 46 & \textbf{68} & 56 & 56 & 52 & 64 & 48 & \underline{60} & 68 & 57 & 8 \\
\quad \EPCRIT{} & 60 & 72 & 64 & \textbf{60} & \textbf{72} & 52 & \textbf{76} & \textbf{56} & 60 & \underline{72} & 64 & 8 \\
\quad \SEMCRIT{} & 40 & 52 & \underline{68} & 44 & \underline{64} & \textbf{64} & 44 & \underline{56} & 60 & \textbf{80} & 57 & 12 \\
\quad \EPSEMCRIT{} & \underline{64} & \underline{76} & 68 & 56 & 60 & 44 & \underline{76} & 52 & 56 & 68 & 62 & 10 \\
\bottomrule
\end{tabular}
\caption{Per-user accuracy for Steam Preference.}
\label{tab:per_user_steam}
\end{subtable}
\hfill
\begin{subtable}[t]{0.49\textwidth}
\centering
\scriptsize
\setlength{\tabcolsep}{2.5pt}
\begin{tabular}{lcccccccccccc}
\toprule
\makecell{Model and \\ Experiment} & U.1 & U.2 & U.3 & U.4 & U.5 & U.6 & U.7 & U.8 & U.9 & U.10 & Avg & Std \\
\midrule
\multicolumn{13}{l}{\textbf{gpt-4o-mini-2024-07-18}} \\
\quad \ZEROSHOT{} & \underline{56} & \textbf{80} & 64 & 40 & 60 & 40 & 60 & 48 & 56 & 52 & 56 & 11 \\
\quad \EPLABEL & \textbf{60} & 68 & 60 & \underline{52} & \textbf{64} & 60 & 56 & 52 & 52 & 56 & 58 & 5 \\
\quad \EPCRIT{} & 56 & 72 & \underline{72} & \textbf{60} & 56 & \textbf{64} & \underline{68} & \underline{56} & \textbf{72} & \underline{72} & 65 & 7 \\
\quad \SEMCRIT{} & 52 & \underline{76} & 68 & 36 & \underline{64} & 56 & 68 & 48 & 56 & 64 & 59 & 11 \\
\quad \EPSEMCRIT{} & 56 & 72 & \textbf{80} & 44 & 60 & \underline{64} & \textbf{76} & \textbf{64} & \underline{72} & \textbf{80} & 67 & 11 \\
\hline
\multicolumn{13}{l}{\textbf{gpt-5-2025-08-07}} \\
\quad \ZEROSHOT{} & 48 & 64 & 56 & 56 & 52 & 32 & 56 & 52 & 52 & 52 & 52 & 8 \\
\quad \EPLABEL & \textbf{76} & \textbf{68} & \underline{80} & 68 & \underline{64} & 64 & \textbf{72} & \textbf{72} & \underline{72} & 76 & 71 & 5 \\
\quad \EPCRIT{} & 64 & \underline{68} & 64 & \textbf{76} & \textbf{68} & 64 & \underline{72} & 64 & \textbf{76} & \textbf{84} & 70 & 7 \\
\quad \SEMCRIT{} & 68 & 68 & 80 & 64 & 64 & \underline{68} & 60 & \underline{72} & 64 & 48 & 66 & 8 \\
\quad \EPSEMCRIT{} & \underline{76} & 68 & \textbf{88} & \underline{76} & 64 & \textbf{76} & 72 & 64 & 68 & \underline{80} & 73 & 7 \\
\hline
\multicolumn{13}{l}{\textbf{gpt-oss-20b}} \\
\quad \ZEROSHOT{} & \textbf{60} & 64 & 56 & \underline{36} & 52 & \underline{64} & \underline{64} & 52 & 52 & \textbf{64} & 56 & 8 \\
\quad \EPLABEL & \underline{56} & 68 & \textbf{80} & 28 & 36 & \textbf{72} & 52 & 52 & 56 & 56 & 56 & 15 \\
\quad \EPCRIT{} & 56 & \underline{72} & \underline{72} & 28 & 48 & 60 & 56 & \textbf{60} & \textbf{76} & 56 & 58 & 13 \\
\quad \SEMCRIT{} & 56 & \textbf{80} & 64 & \textbf{44} & \textbf{60} & 60 & 56 & \underline{56} & \underline{72} & 52 & 60 & 10 \\
\quad \EPSEMCRIT{} & 52 & 72 & 72 & 36 & \underline{56} & 56 & \textbf{68} & 52 & 72 & \underline{60} & 60 & 11 \\
\hline
\multicolumn{13}{l}{\textbf{llama-4-scout}} \\
\quad \ZEROSHOT{} & \textbf{56} & 72 & \textbf{68} & \textbf{44} & 52 & 52 & \underline{60} & 52 & 48 & \textbf{64} & 57 & 9 \\
\quad \EPLABEL & 48 & \underline{84} & 48 & 28 & 52 & 52 & 52 & \textbf{64} & 52 & \underline{64} & 54 & 14 \\
\quad \EPCRIT{} & 52 & 76 & \underline{68} & \underline{44} & \textbf{64} & \underline{56} & 60 & \underline{60} & \underline{68} & 60 & 61 & 9 \\
\quad \SEMCRIT{} & \underline{56} & 84 & 68 & 44 & \underline{60} & 44 & \textbf{68} & 44 & 68 & 56 & 59 & 12 \\
\quad \EPSEMCRIT{} & 56 & \textbf{96} & 56 & 44 & 60 & \textbf{60} & 60 & 60 & \textbf{76} & 48 & 62 & 14 \\
\hline
\multicolumn{13}{l}{\textbf{qwen3-235b-a22b-instruct-2507}} \\
\quad \ZEROSHOT{} & 52 & 68 & 56 & 44 & 60 & 40 & \textbf{76} & 52 & 52 & 48 & 55 & 10 \\
\quad \EPLABEL & 60 & 64 & 56 & 52 & 56 & \underline{60} & \underline{68} & 40 & 64 & \textbf{68} & 59 & 8 \\
\quad \EPCRIT{} & 60 & 64 & 68 & \underline{60} & 60 & \textbf{68} & 68 & \textbf{60} & \textbf{72} & 56 & 64 & 5 \\
\quad \SEMCRIT{} & \textbf{68} & \underline{72} & \textbf{80} & 52 & \underline{68} & 36 & 60 & \underline{60} & 48 & 60 & 60 & 12 \\
\quad \EPSEMCRIT{} & \underline{64} & \textbf{76} & \underline{76} & \textbf{68} & \textbf{76} & 56 & 64 & 52 & \underline{68} & \underline{64} & 66 & 8 \\
\hline
\multicolumn{13}{l}{\textbf{qwen3-vl-235b-a22b-thinking}} \\
\quad \ZEROSHOT{} & 44 & 72 & 48 & 32 & \underline{64} & 56 & 52 & \textbf{64} & 56 & 48 & 54 & 11 \\
\quad \EPLABEL & \textbf{64} & 72 & \underline{80} & 48 & 52 & \textbf{72} & 52 & 36 & \textbf{72} & \textbf{72} & 62 & 13 \\
\quad \EPCRIT{} & 52 & \textbf{76} & 64 & 44 & 56 & \underline{60} & \textbf{68} & 44 & 68 & 44 & 58 & 11 \\
\quad \SEMCRIT{} & 40 & 72 & \textbf{84} & \textbf{60} & \textbf{76} & 48 & \underline{64} & 36 & 60 & 48 & 59 & 15 \\
\quad \EPSEMCRIT{} & \underline{60} & \underline{76} & 76 & \underline{56} & 56 & 60 & 60 & \underline{48} & \underline{72} & \underline{56} & 62 & 9 \\
\bottomrule
\end{tabular}
\caption{Per-user accuracy for Book Preference.}
\label{tab:per_user_book}
\end{subtable}
\caption{For each model and user, the highest score is \textbf{bolded} and the second-highest is \underline{underlined}.}
\end{table*}

\begin{table*}[ht]
\centering
\begin{subtable}[t]{0.49\textwidth}
\centering
\scriptsize
\setlength{\tabcolsep}{2.5pt}
\begin{tabular}{lcccccccccccc}
\toprule
\makecell{Model and \\ Experiment} & U.1 & U.2 & U.3 & U.4 & U.5 & U.6 & U.7 & U.8 & U.9 & U.10 & Avg & Std \\
\midrule
\multicolumn{13}{l}{\textbf{gpt-4o-mini-2024-07-18}} \\
\quad \ZEROSHOT{} & 32 & 72 & 44 & \textbf{68} & \underline{64} & \underline{52} & \underline{68} & 52 & \textbf{56} & 56 & 56 & 12 \\
\quad \EPLABEL & \textbf{48} & 72 & \textbf{68} & 56 & 64 & 36 & 68 & 60 & \underline{48} & 56 & 58 & 11 \\
\quad \EPCRIT{} & \underline{44} & \underline{76} & \underline{68} & \underline{64} & 64 & \textbf{56} & 64 & \underline{72} & 44 & \underline{68} & 62 & 10 \\
\quad \SEMCRIT{} & 36 & 64 & 48 & 48 & 64 & 48 & 64 & 60 & 36 & 60 & 53 & 11 \\
\quad \EPSEMCRIT{} & 32 & \textbf{80} & 64 & 64 & \textbf{72} & 36 & \textbf{76} & \textbf{80} & 36 & \textbf{76} & 62 & 18 \\
\hline
\multicolumn{13}{l}{\textbf{gpt-5-2025-08-07}} \\
\quad \ZEROSHOT{} & 44 & \underline{76} & 64 & 36 & 64 & 44 & \textbf{72} & 52 & 52 & 52 & 56 & 12 \\
\quad \EPLABEL & \underline{68} & \textbf{84} & \textbf{96} & \underline{80} & 60 & \textbf{64} & \underline{64} & 88 & 52 & 64 & 72 & 13 \\
\quad \EPCRIT{} & \textbf{72} & 68 & \underline{80} & \textbf{92} & 64 & \underline{64} & 64 & \textbf{92} & \underline{60} & \textbf{76} & 73 & 11 \\
\quad \SEMCRIT{} & 56 & 48 & 68 & 60 & \underline{68} & 52 & 48 & 56 & \textbf{68} & 56 & 58 & 7 \\
\quad \EPSEMCRIT{} & 60 & 68 & 68 & 76 & \textbf{72} & 60 & 64 & \underline{92} & 52 & \underline{76} & 69 & 11 \\
\hline
\multicolumn{13}{l}{\textbf{gpt-oss-20b}} \\
\quad \ZEROSHOT{} & 24 & 60 & 52 & 52 & \textbf{68} & 44 & 52 & 64 & \underline{56} & 56 & 53 & 12 \\
\quad \EPLABEL & 32 & 68 & 52 & 52 & 44 & \textbf{48} & 52 & 72 & 48 & 68 & 54 & 12 \\
\quad \EPCRIT{} & \textbf{60} & 72 & \underline{72} & 48 & \underline{60} & \underline{48} & \underline{60} & \textbf{80} & 44 & 64 & 61 & 11 \\
\quad \SEMCRIT{} & \underline{52} & \textbf{76} & 52 & \underline{56} & 60 & 48 & \textbf{84} & 64 & 36 & \textbf{76} & 60 & 14 \\
\quad \EPSEMCRIT{} & 52 & \underline{76} & \textbf{76} & \textbf{60} & 48 & 44 & 44 & \underline{76} & \textbf{64} & \underline{72} & 61 & 13 \\
\hline
\multicolumn{13}{l}{\textbf{llama-4-scout}} \\
\quad \ZEROSHOT{} & \textbf{40} & 56 & 48 & \textbf{60} & 56 & 44 & 56 & 56 & \textbf{52} & 72 & 54 & 8 \\
\quad \EPLABEL & \underline{36} & \underline{72} & \underline{52} & 56 & \textbf{76} & \textbf{56} & 60 & \textbf{72} & 48 & \textbf{76} & 60 & 13 \\
\quad \EPCRIT{} & 36 & 64 & \textbf{64} & \underline{60} & \underline{64} & 44 & 60 & 64 & 48 & \underline{76} & 58 & 11 \\
\quad \SEMCRIT{} & 28 & \textbf{76} & 36 & 60 & 64 & \underline{48} & \textbf{64} & 68 & \underline{52} & 72 & 57 & 15 \\
\quad \EPSEMCRIT{} & 28 & 68 & 52 & 60 & 64 & 48 & \underline{64} & \underline{72} & 52 & 68 & 58 & 12 \\
\hline
\multicolumn{13}{l}{\textbf{qwen3-235b-a22b-instruct-2507}} \\
\quad \ZEROSHOT{} & 52 & 60 & 44 & 32 & 56 & 48 & \underline{64} & 60 & \textbf{64} & 48 & 53 & 10 \\
\quad \EPLABEL & \underline{64} & 64 & 76 & \underline{52} & 56 & \textbf{60} & 56 & \underline{84} & 60 & \underline{76} & 65 & 10 \\
\quad \EPCRIT{} & \textbf{72} & \underline{68} & \underline{84} & \textbf{64} & \textbf{64} & 44 & 56 & \textbf{92} & \underline{64} & \textbf{88} & 70 & 14 \\
\quad \SEMCRIT{} & 48 & 68 & 68 & 44 & \underline{60} & 52 & \textbf{68} & 72 & 64 & 72 & 62 & 10 \\
\quad \EPSEMCRIT{} & 64 & \textbf{72} & \textbf{88} & 52 & 60 & \underline{60} & 60 & 68 & 60 & 68 & 65 & 9 \\
\hline
\multicolumn{13}{l}{\textbf{qwen3-vl-235b-a22b-thinking}} \\
\quad \ZEROSHOT{} & 40 & \textbf{68} & 44 & 52 & \textbf{76} & 36 & \textbf{68} & 60 & \underline{56} & 56 & 56 & 12 \\
\quad \EPLABEL & \textbf{56} & \underline{68} & 64 & \textbf{80} & \underline{72} & 44 & \underline{68} & 80 & 52 & 64 & 65 & 11 \\
\quad \EPCRIT{} & \underline{56} & 64 & \textbf{72} & \underline{68} & 60 & \underline{48} & 52 & \textbf{92} & \textbf{60} & 64 & 64 & 12 \\
\quad \SEMCRIT{} & 56 & 64 & \underline{68} & 40 & 52 & 40 & 52 & 56 & 48 & \textbf{68} & 54 & 10 \\
\quad \EPSEMCRIT{} & 52 & 56 & 68 & 60 & 64 & \textbf{68} & 60 & \underline{84} & 56 & \underline{68} & 64 & 9 \\
\bottomrule
\end{tabular}
\caption{Per-user accuracy for Anime Preference.}
\label{tab:per_user_anime}
\end{subtable}
\begin{subtable}[t]{0.49\textwidth}
\centering
\scriptsize
\setlength{\tabcolsep}{2.5pt}
\begin{tabular}{lcccccccccccc}
\toprule
\makecell{Model and \\ Experiment} & U.1 & U.2 & U.3 & U.4 & U.5 & U.6 & U.7 & U.8 & U.9 & U.10 & Avg & Std \\
\midrule
\multicolumn{13}{l}{\textbf{gpt-4o-mini-2024-07-18}} \\
\quad \ZEROSHOT{} & 60 & 36 & 72 & \textbf{56} & 20 & \underline{68} & \underline{36} & 60 & \textbf{64} & 44 & 52 & 16 \\
\quad \EPLABEL & 36 & 52 & 92 & 24 & 24 & 20 & \textbf{48} & 64 & 48 & \textbf{48} & 46 & 21 \\
\quad \EPCRIT{} & 60 & \underline{64} & \underline{96} & \underline{52} & \textbf{40} & 68 & 32 & \textbf{72} & \underline{60} & 40 & 58 & 18 \\
\quad \SEMCRIT{} & \textbf{76} & 52 & 88 & 52 & 16 & \textbf{72} & 28 & 60 & 60 & \underline{48} & 55 & 20 \\
\quad \EPSEMCRIT{} & \underline{64} & \textbf{68} & \textbf{100} & 52 & \underline{40} & 64 & 36 & \underline{68} & 56 & 36 & 58 & 18 \\
\hline
\multicolumn{13}{l}{\textbf{gpt-5-2025-08-07}} \\
\quad \ZEROSHOT{} & 60 & 32 & 72 & 44 & 20 & 40 & \textbf{56} & \textbf{84} & 40 & 40 & 49 & 18 \\
\quad \EPLABEL & \textbf{64} & \underline{60} & 88 & \textbf{88} & \textbf{56} & \textbf{76} & \underline{52} & 72 & \underline{76} & 60 & 69 & 12 \\
\quad \EPCRIT{} & 56 & 56 & 88 & \underline{88} & 44 & \underline{72} & 44 & 72 & \textbf{88} & 48 & 66 & 17 \\
\quad \SEMCRIT{} & \underline{64} & 28 & \textbf{92} & 84 & 52 & 60 & 52 & \underline{80} & 44 & \textbf{80} & 64 & 19 \\
\quad \EPSEMCRIT{} & 60 & \textbf{64} & \underline{92} & 84 & \underline{56} & 72 & 32 & 72 & 72 & \underline{72} & 68 & 16 \\
\hline
\multicolumn{13}{l}{\textbf{gpt-oss-20b}} \\
\quad \ZEROSHOT{} & \underline{56} & \underline{48} & 76 & 36 & 28 & 60 & 44 & \textbf{76} & 56 & 40 & 52 & 15 \\
\quad \EPLABEL & 40 & 48 & \textbf{84} & 60 & 44 & 20 & 44 & \underline{76} & \underline{72} & \underline{52} & 54 & 18 \\
\quad \EPCRIT{} & 52 & \textbf{52} & 80 & \textbf{80} & \underline{52} & \textbf{72} & 48 & 72 & \textbf{80} & 52 & 64 & 13 \\
\quad \SEMCRIT{} & \textbf{64} & 32 & \underline{84} & 36 & 32 & 64 & \textbf{56} & 68 & 40 & 48 & 52 & 17 \\
\quad \EPSEMCRIT{} & 56 & 44 & 84 & \underline{72} & \textbf{60} & \underline{72} & \underline{52} & 68 & 68 & \textbf{60} & 64 & 11 \\
\hline
\multicolumn{13}{l}{\textbf{llama-4-scout}} \\
\quad \ZEROSHOT{} & 52 & \underline{40} & 68 & \textbf{48} & \textbf{32} & \textbf{76} & 40 & 60 & 44 & 36 & 50 & 14 \\
\quad \EPLABEL & \underline{56} & 40 & \textbf{96} & \underline{40} & 24 & 64 & 48 & 72 & \underline{68} & 40 & 55 & 20 \\
\quad \EPCRIT{} & \textbf{64} & 36 & 88 & 40 & \underline{32} & \underline{68} & \underline{52} & 72 & \textbf{76} & \textbf{48} & 58 & 18 \\
\quad \SEMCRIT{} & 44 & 32 & 92 & 36 & 20 & 64 & 48 & \underline{76} & 48 & 44 & 50 & 20 \\
\quad \EPSEMCRIT{} & 56 & \textbf{48} & \underline{96} & 36 & 20 & 64 & \textbf{56} & \textbf{80} & 60 & \underline{48} & 56 & 20 \\
\hline
\multicolumn{13}{l}{\textbf{qwen3-235b-a22b-instruct-2507}} \\
\quad \ZEROSHOT{} & 56 & 52 & 60 & 40 & 20 & \underline{64} & \underline{48} & 52 & \textbf{68} & \underline{56} & 52 & 13 \\
\quad \EPLABEL & 44 & 64 & 84 & 40 & 48 & 36 & 36 & 56 & 56 & 56 & 52 & 14 \\
\quad \EPCRIT{} & 56 & 48 & \underline{92} & \underline{64} & 60 & \textbf{68} & 24 & 64 & 60 & 48 & 58 & 16 \\
\quad \SEMCRIT{} & \textbf{60} & \textbf{72} & 92 & 56 & \underline{72} & 56 & \textbf{52} & \underline{68} & 40 & \textbf{60} & 63 & 13 \\
\quad \EPSEMCRIT{} & \underline{60} & \underline{72} & \textbf{96} & \textbf{68} & \textbf{76} & 56 & 32 & \textbf{72} & \underline{64} & 44 & 64 & 17 \\
\hline
\multicolumn{13}{l}{\textbf{qwen3-vl-235b-a22b-thinking}} \\
\quad \ZEROSHOT{} & 52 & 44 & 80 & 40 & 20 & 64 & 44 & 60 & 52 & 36 & 49 & 16 \\
\quad \EPLABEL & \underline{56} & \textbf{64} & 88 & \textbf{84} & \textbf{64} & 64 & 40 & 60 & 60 & 36 & 62 & 15 \\
\quad \EPCRIT{} & \textbf{60} & \underline{60} & 88 & \underline{84} & \underline{60} & \underline{72} & \textbf{48} & \textbf{80} & \textbf{76} & 44 & 67 & 14 \\
\quad \SEMCRIT{} & 52 & 48 & \underline{92} & 80 & 48 & \textbf{84} & 40 & 56 & 68 & \textbf{52} & 62 & 17 \\
\quad \EPSEMCRIT{} & 56 & 48 & \textbf{96} & 80 & 48 & 68 & \underline{48} & \underline{72} & \underline{76} & \underline{48} & 64 & 16 \\
\bottomrule
\end{tabular}
\caption{Per-user accuracy for Movie Preference.}
\label{tab:per_user_movie}
\end{subtable}
\caption{For each model and user, the highest score is \textbf{bolded} and the second-highest is \underline{underlined}.}
\end{table*}

